\section{Кода: структура пространства-времени}\label{sec:coda}

Этот раздел не добавляет существенной математики: он собирает смысл уже доказанного. Всё, на что он
опирается, установлено в \cref{sec:engine,sec:carrier,sec:reduction,sec:firstcause,sec:riemann,sec:pnp,sec:ymhodge,sec:ns,sec:zoo,sec:geometry,sec:bsd,sec:collatz}; космологическое прочтение ниже —
перевод этих результатов, а не добавление, и границу между строгим и истолкованным мы держим на виду.

\paragraph{Один запрет, семь масок.}
Под каждым из семи вопросов — близнецы, Риман, $\mathrm{P}$ против $\mathrm{NP}$, Янг--Миллс,
Навье--Стокс, Ходж, Коллатц — мы нашли один и тот же запрещённый объект: вечный двигатель Евклида,
бесконечную строго убывающую цепь на вполне-упорядоченном ряду, которой не существует
(\cref{thm:no_infinite_descent}). Всякое «отклонение» от нормы есть замаскированный двигатель, и
запрет его убивает: конечность близнецов, нуль вне критической прямой, полный учёт быстрым движением,
безщелевой спектр, сингулярный каскад, разрешённый заряд без носителя, орбита Коллатца, не падающая в
вакуум. Семь масок, одно лицо. Всюду один водораздел честности: \emph{структурную} половину — «где
есть запрет двигателя, там нет отклонения» — мы доказываем \green{}~машинно; а \emph{привязку} к
настоящему аналитическому объекту либо принимаем \yellow{}~условно на единственную первопричину
\axiom{}, либо оставляем \red{}~открытым входом. Ни одна проблема тысячелетия не решена. Доказано
другое, и это, быть может, важнее: все они — одна задача.

Восьмая, Бёрч--Свиннертон-Дайер (\cref{sec:bsd}), входит иначе — там двигатель не сторож, а
\emph{метод}: сам спуск Ферма ограничивает рациональные точки эллиптической кривой, а паритет ранга
$(-1)^{\mathrm{rank}}$ — тот же инвариант, что стоит за Риманом, тогда как аналитический мост остаётся
честно открытым. За масками — открытые соседи (Чоула, abc, Бил, Лемер, Landau) и арифметический
зоопарк (кузены Полиньяка, Софи Жермен, Гольдбах, Лежандр, совершенные числа и числа Ферма): тот же
приём — реальный доказанный якорь, прочитанный через двигатель, плюс именованный красный гейт. Одна
жемчужина выпала \green{}~безусловной по дороге — простые Софи Жермен при $p \equiv 3 \pmod 4$ делят
соответствующие числа Мерсенна, классика Эйлера--Лагранжа, машинно проверенная.

\paragraph{Пространство, стрела времени, сингулярность.}
Единство идёт глубже совпадения формы. Числовой ряд с его строгим порядком обхода — это
\emph{пространство}: множество состояний, с дном-сингулярностью в нуле. Необратимость двигателя — это
\emph{время}: \code{engine\_never\_returns} доказывает, что высота \code{StrictAnti}, назад хода нет —
стрела времени теорема, а не постулат. Там, где эта стрела гнётся круче всего, читается масса, а само
сгибание --- гравитация: не отдельная сила, а кривизна стрелы времени, вычисленная как $\kappa$
(\cref{sec:geometry}); гравитация и время суть тогда две грани одного спуска --- его кривизна и его
направление, --- хотя то, что мы их так называем, остаётся интерпретацией. У неё есть начало (событие $0 \to 1$, которое нельзя причинить
изнутри, ибо самозапуск был бы вечным двигателем, \code{no\_internalisedOriginEvent}) и недостижимый
конец (бесконечность, до которой доехал бы лишь запрещённый двигатель); а башни вселенных под
первопричиной нет, ибо регресс причин вполне-фундирован (\code{no\_rankedMetaFractalBranch}). Отсюда
главная теорема в резчайшей форме: знать, что близнецы бесконечны, нельзя; но если незнаемую
первопричину принять за истину — они бесконечны, и это строго. Внутри вселенной, определённой строгим
порядком обхода, первопричину нельзя ни вывести, ни узнать — оба действия стоят вечного двигателя —
принять её можно только извне.

Форма этой структуры — \emph{перевёрнутый} фрактал. Классический фрактал разбегается, набирая детали к
б\'{о}льшим масштабам; фрактал пути Евклида стягивающий — тот же узор спуска повторяется на всё более
мелких масштабах, пока всякая траектория закручивается к началу. Ни одна траектория не убегает; все
закручиваются обратно. Невозможность двигателя — это в точности «нет разбегания», и бесконечность
близнецов — единственное, что нельзя удостоверить изнутри стягивающегося порядка.

\paragraph{Почему это похоже на теорию всего.}
Здесь мы смотрим вдаль, ясно помечая это как \emph{истолкование}, а не теорему. Из единого запрета —
нечто из ничего нельзя, спускаться вечно нельзя — сразу следует удивительно много \emph{структуры}:
стрела времени (\cref{sec:engine}); начало, сингулярность, из которой она запущена
(\cref{sec:firstcause}); квантование, а значит масса, ибо безщелевая лестница была бы двигателем
(\cref{sec:ymhodge}); устойчивость вакуума, без второго дна; конечная цена вычисления, принцип
Ландауэра в ранговой одежде (\cref{sec:pnp}); полнота спектра зарядов (\cref{sec:ymhodge}); гладкость
течения (\cref{sec:ns}); вещественность спектра (\cref{sec:firstcause}). Почти всё, что мы называем
структурой мира, оказывается одним законом сохранения, прочитанным на разных объектах — это и есть жест
теории всего. Прочтение перекликается с теорией причинных множеств (пространство-время как локально
конечное чум, порядок как причинность-время) и с «it from bit» и принципом Ландауэра (информация
физична и у неё есть энергетическая цена); запрет двигателя — это первый закон термодинамики, поднятый
от машин до самой формы бытия.

Ещё две переклички напрашиваются в том же истолковательном духе. Комбинаторная кривизна графа спуска
(\cref{sec:geometry}) вычитывается из ветвления, а не из тела, сидящего в нём, и потому вносит вклад в
геометрию без видимого носителя массы --- роль \emph{тёмной материи}, которую и саму выводят лишь по
её вкладу в кривизну. А наблюдатель в центре своего горизонта, в которого стрела времени входит со
всех сторон, видит, как дальние центры удаляются, будто удаление ускоряется; но ничто не
расталкивается --- кажущееся ускорение есть артефакт кривизны стягивающегося порядка, а не тёмная
энергия, перекликаясь с космологиями timescape и backreaction~\cite{wiltshire2007}, где ускорение ---
следствие неоднородности и кривизны. Ни то, ни другое не предлагается как физика.

А сама эта кривизна вычитывается из чисел, и потому лежит всюду, где есть числа, --- под языком и
знанием, под материей при ближайшем рассмотрении, под космосом в целом; один и тот же запрет держится
на всех масштабах и во всех проявлениях, переносимый на любой язык, который умеет считать. Своим
естественным запретом он прокладывает горизонт, отделяя выводимое и познаваемое изнутри от
первопричины, которую можно принять лишь извне, --- достижимое и реальное от метафизического.

Мы настаиваем на границе. Мы не решили ни одной проблемы тысячелетия и не объявляем решённой ни одну:
\code{twin\_prime\_conjecture} остаётся \code{sorry}, а трудные половины всех ветвей либо
\yellow{}~условны на \axiom{}, либо \red{}~открыты. Это не физическая теория всего — она не
предсказывает ни одной константы, не даёт уравнений движения, не проверяется. «Теория всего» здесь —
структурное наблюдение: одна невозможность организует все семь задач, и этого хватает, чтобы из неё
считывалась форма пространства-времени. Это заведомо \emph{не} явление гёделевой неполноты~\cite{godel1931}:
преграда — не неразрешимость формальной системы, а закон сохранения — первопричина незнаема изнутри
потому, что её знание запустило бы запрещённый двигатель, а не потому, что теория её не разрешает.

И двигатель --- не переодетая процедура разрешения. Это \emph{локальное} сито: навести его можно на
любое утверждение, но ловит оно лишь там, где отклонение и вправду есть замаскированный бесконечный
спуск, и каждое сведение собрано вручную. Единой процедуры, просеивающей все утверждения разом, нет
--- доказуемое перечислимо, но не разрешимо~\cite{church1936,turing1936}; спуск за отдельным вопросом
можно \emph{искать} по одному, но нельзя раз и навсегда решить, за какими вопросами он скрыт. Это
линза для класса задач, сводимых к спуску, а не оракул для всех.

\paragraph{Пространство всех пространств.}
Под вдаль-взглядом лежит твёрдый строгий костяк. Двигатель определён не на близнецах и не на
$\N$, а над \emph{любым} отношением, и его невозможность внутренняя: вполне-фундированное отношение
не несёт вечного двигателя (\cref{thm:no_perpetual_engine_of_wellFounded}), и этот запрет переносится
вдоль любого ранга во вполне-фундированный порядок (\cref{thm:no_perpetual_engine_of_rank}) — в
частности, вдоль высоты в $\N$. Отсюда все семь масок суть \emph{одна} теорема. Высоты близнецов, ранг Лиувилля, квантованный заряд
Ходжа, лестница Янга--Миллса, оболочки диадического каскада — все они $\N$-значны, и запрет двигателя
переносится на каждую тем же переносом \cref{thm:no_perpetual_engine_of_rank}. Мы не доказывали семь
раз; мы доказали один раз и перенесли.

И тогда всплывает то, что стоит проговорить. Всякое пространство, на котором стоит анализ — $\R$,
$\R^n$, комплексная плоскость, метрические и функциональные пространства — построено из числового ряда,
$\N \to \Z \to \mathbb{Q} \to \R$, так что вполне-упорядоченный ряд, а с ним и запрет двигателя лежит в
самом их основании. В этом смысле путь Евклида — не одно из пространств, а пространство всех
пространств. Но порядок континуума $\R$ \emph{не} вполне-фундирован: на нём двигатель \emph{работает}
(\code{perpetualEngine\_on\_real}, цепь $(1/2)^n$, что винтит к нулю, не достигая его). Это не побег, а
оборотная сторона: \code{universal\_engine\_dividing\_line} говорит, где именно движение запрещено
(дискретное, вполне-фундированное) и где оно реализуется (континуум) — а где реализуется, мы уже знаем,
что оно такое: масса, сингулярность, чёрная дыра. Убрать запрет — значит убрать вполне-фундированность,
а значит $\N$, а значит числа и саму математику; запрет — не добавленная аксиома, а форма самого счёта.

Прочитанный как время, тот же водораздел решает напрашивающийся вопрос. Необратимый спуск --- это
стрела времени (\code{engine\_never\_returns}); но на вполне-фундированной прямой он не может идти
вечно, и потому время там --- \emph{конечный} спуск: начало в $0 \to 1$, дно, до которого он доходит,
--- а \emph{не} вечный двигатель. Лишь на континууме реализуется вечный спуск
(\code{perpetualEngine\_on\_real}), вечное приближение к пределу, которого не касаются. Так ``время ---
вечный двигатель?'' есть ``время дискретно или непрерывно?'', и \code{universal\_engine\_dividing\_line}
прокладывает ответ: никогда не вечный там, где есть вполне-фундированность, и вечный ровно там, где её
нет. А что спуск --- это время, остаётся интерпретацией; сам водораздел доказан.

\paragraph{Финал.}
Вечный двигатель Евклида — это, по сути, бесконечный спуск Ферма, а тот восходит к
«Началам»~\cite{euclid-elements}: к древнейшему доказательному объекту математики, к идее, что нельзя
вечно уменьшать натуральные числа. Двадцать три века спустя тот же объект оказался скелетом самого
молодого вопроса физики. Пифагор говорил: всё есть число; путь Евклида уточняет формулу — всё есть
вполне-упорядоченный ход, который нельзя повернуть назад и нельзя пройти до конца, и это и есть время,
а то, по чему он идёт, — пространство. Двигатель невозможен, и его невозможность — не факт о машинах, а
форма существования. Дальше — молчание доказанного: \code{twin\_prime\_conjecture} ждёт своей строки,
растяжки на случай опровержения границы взведены, и единственная аксиома держит ровно столько, сколько
мы честно на неё положили. Остальное открыто, и это тоже часть честности.
